\chapter{Guided Lab and Exercise Compendium}

The weekly chapters explain what students are learning. This compendium gives the dsDNA Core program a deeper set of labs and real-world exercises to fill full-day work blocks. Each lab ends with a saved artifact, a research log entry, and a short review check.

\section{Lab 1: Chemical Question Triage}
\textbf{Objective:} Teach students to reject questions that cannot be answered with available chemical evidence.

\textbf{Inputs:} Three candidate research interests, the project brief worksheet, and a program-provided example chemical list.

\textbf{Procedure:}
\begin{enumerate}
  \item For each candidate question, identify the system, comparison, chemical evidence type, expected source tables, and likely limitation.
  \item Mark each question as feasible, feasible with narrower scope, or not feasible in 10 weeks.
  \item Rewrite weak questions so they specify a chemical set, comparison, evidence type, and output.
  \item Choose one question for the internship and one backup question.
\end{enumerate}

\textbf{Quality gate:} The review should make clear what data will be analyzed and what decision the analysis could support.

\textbf{Real-world exercise:} A collaborator asks whether these GC-MS hits can prove that a plant treatment increases medicinal compounds. Students must rewrite the request into a defensible question that separates tentative identity, trait enrichment, and follow-up validation.

\section{Lab 2: Reproducible Project Build}
\textbf{Objective:} Build a project that can be reopened and rerun by another person.

\textbf{Inputs:} RStudio, Git, uafR source or installed package, and \path{training/scripts/01_project_setup.R}.

\textbf{Procedure:}
\begin{enumerate}
  \item Run \texttt{create\_training\_project()}.
  \item Add a project title, research question, data description, and script order to \path{README.md}.
  \item Create \path{logs/session_info.txt} using \texttt{capture.output(sessionInfo())}.
  \item Commit only appropriate files. Confirm that tokens, private data, and local machine paths are absent.
\end{enumerate}

\textbf{Quality gate:} The project runs from its root directory and contains a log explaining what has been done so far.

\textbf{Real-world exercise:} Students exchange project folders and try to identify the correct script order without asking the author. Any confusion becomes a README revision.

\section{Lab 3: Database Source Audit}
\textbf{Objective:} Teach students to interpret database sources as evidence with scope and limitations.

\textbf{Inputs:} Three chemicals from the student's project, PubChem pages or uafR PubChem outputs, KEGG records when available, and the source comparison worksheet.

\textbf{Procedure:}
\begin{enumerate}
  \item For each chemical, identify a stable name, CID if available, and likely ambiguity.
  \item Record which sources provide identity, properties, hazard, occurrence, sensory, biomedical, pathway, reaction, and literature evidence.
  \item Mark every source-backed field as analysis-ready, narrative-only, or not relevant.
  \item Write one sentence that states what each source can support and one sentence that states what it cannot support.
\end{enumerate}

\textbf{Quality gate:} Students must avoid treating database presence as proof of biological presence in their sample.

\textbf{Real-world exercise:} Students compare caffeine and methyl salicylate. They should notice that one chemical may have strong biomedical and food-context evidence, while another may connect strongly to plant and sensory contexts. The result is a source suitability comparison, not a claim that one chemical is more important.

\section{Lab 4: Offline Core Workflow}
\textbf{Objective:} Teach the GC-MS workflow without depending on live web services.

\textbf{Inputs:} \path{training/scripts/04_core_workflow.R}, bundled \texttt{standard\_spread}, and bundled \texttt{standard\_exacto}.

\textbf{Procedure:}
\begin{enumerate}
  \item Source the script and run \texttt{run\_core\_workflow(live\_lookup = FALSE)}.
  \item Inspect the structure of \texttt{spread} and \texttt{exact}.
  \item Identify which columns would be needed to write a methods paragraph.
  \item If services are available and live requests are appropriate, rerun with \texttt{live\_lookup = TRUE} and compare outputs.
\end{enumerate}

\textbf{Quality gate:} Students can explain which objects are bundled teaching outputs and which require live lookup.

\textbf{Real-world exercise:} A student receives a CSV with the wrong column names. They must create an input audit that lists missing columns, suspicious duplicate keys, and a decision on whether the file can be processed.

\section{Lab 5: Research Enrichment Smoke Test}
\textbf{Objective:} Run a small \texttt{categorate()} workflow and inspect validation before interpretation.

\textbf{Inputs:} \path{training/scripts/05_categorate_research.R}, two or three common compounds, and \texttt{library\_data}.

\textbf{Procedure:}
\begin{enumerate}
  \item Run a two-compound test with \texttt{assay\_detail\_limit = 0}.
  \item Inspect \texttt{ValidationSummary}, \texttt{TableQuality}, and \texttt{SourceCoverage}.
  \item Save the result as an RDS file and write the query settings in the research log.
  \item Identify three useful tables and three tables that are empty or weak for the chosen compounds.
\end{enumerate}

\textbf{Quality gate:} Students do not inspect \texttt{ChemicalTraitMatrix} until they have inspected validation and source coverage.

\textbf{Real-world exercise:} The student must decide whether to run 50 compounds immediately or batch-test 5 first. They should justify the decision using database etiquette, caching, and source coverage.

\section{Lab 6: Evidence Court}
\textbf{Objective:} Train students to defend or withdraw claims based on evidence rows.

\textbf{Inputs:} \texttt{ChemicalTraitMatrix}, \texttt{ChemicalTraitEvidence}, \texttt{ChemicalTraitReport}, and one draft result claim.

\textbf{Procedure:}
\begin{enumerate}
  \item Choose one matrix column that appears important.
  \item Pull all evidence rows for that matrix key.
  \item Identify source database, source field, evidence text, source URL, confidence, and extraction rule.
  \item Decide whether the trait can support a result, a hypothesis, or only a note.
  \item Rewrite the claim using the correct strength of evidence.
\end{enumerate}

\textbf{Quality gate:} Every accepted claim must trace to evidence and every rejected claim must be documented.

\textbf{Real-world exercise:} Students hold a mock review panel. One student argues for a claim, one argues against it, and one acts as editor. The editor decides whether the claim belongs in Results, Discussion, or Future Work.

\section{Lab 7: Trait Matrix Design}
\textbf{Objective:} Build matrices that answer a research question rather than maximize column count.

\textbf{Inputs:} \texttt{ChemicalTraits}, \texttt{ChemicalTraitOntology}, \texttt{chemicalTraitMatrix()}, and the grouping dictionary worksheet.

\textbf{Procedure:}
\begin{enumerate}
  \item Build a core binary matrix with medium confidence or higher.
  \item Build a wider confidence-weighted matrix.
  \item Compare matrix dimensions, sparsity, and interpretability.
  \item Select a final matrix profile and document why weaker columns were removed.
\end{enumerate}

\textbf{Quality gate:} The final matrix should be smaller, clearer, and more defensible than the widest possible matrix.

\textbf{Real-world exercise:} A project review requests clustering. Students must decide whether the matrix has enough meaningful columns and coverage to make clustering interpretable.

\section{Lab 8: Figure Clinic}
\textbf{Objective:} Convert validated matrices and summaries into figures that support interpretation.

\textbf{Inputs:} \path{training/scripts/07_visualization.R}, a saved enrichment result, and the figure planning worksheet.

\textbf{Procedure:}
\begin{enumerate}
  \item Generate a source coverage heatmap and a trait matrix heatmap.
  \item Build one additional figure that matches the project question.
  \item Write a caption that names the source table, evidence type, and limitation.
  \item Peer review the figure using the presentation rubric.
\end{enumerate}

\textbf{Quality gate:} A figure must answer a question, not merely display available columns.

\textbf{Real-world exercise:} Students receive two versions of a figure: one visually attractive but overclaiming, one plain but honest. They revise the attractive figure so it remains readable while becoming scientifically defensible.

\section{Lab 9: Literature-to-Evidence Bridge}
\textbf{Objective:} Connect database-derived patterns to peer-reviewed context.

\textbf{Inputs:} A figure, an evidence row, and three papers related to the system or chemical class.

\textbf{Procedure:}
\begin{enumerate}
  \item For each paper, extract organism or system, method, chemical evidence, main claim, and limitation.
  \item Map each paper to a specific project result.
  \item Mark whether the paper supports background, interpretation, mechanism, or follow-up.
  \item Write a claim-evidence-limitation paragraph.
\end{enumerate}

\textbf{Quality gate:} Literature should sharpen interpretation, not inflate it.

\textbf{Real-world exercise:} A student finds a paper about the same compound but a different organism and analytical method. They must decide whether it belongs in Introduction, Discussion, or not at all.

\section{Lab 10: Final Handoff Audit}
\textbf{Objective:} Ensure the project can survive after the student leaves.

\textbf{Inputs:} Final project folder, \path{training/scripts/09_reproducibility_check.R}, final report, figures, and saved result objects.

\textbf{Procedure:}
\begin{enumerate}
  \item Run \texttt{check\_training\_project()}.
  \item Confirm every figure can be regenerated from scripts.
  \item Confirm every major result has a source table, evidence row, or citation.
  \item Move abandoned scratch files into an archive folder or remove them after program review.
  \item Write a continuation note for the next student.
\end{enumerate}

\textbf{Quality gate:} A reviewer should be able to rerun or inspect the work without asking the student to explain where files are.

\textbf{Real-world exercise:} Students trade final project folders and perform a handoff audit. The reviewer writes a one-page note naming what they could reproduce, what was unclear, and what should be fixed first.
